\documentclass[UTF8,a4paper,12pt]{ctexart}

\usepackage{amsmath,amssymb,bm}
\usepackage{booktabs}
\usepackage{geometry}
\usepackage{longtable}
\usepackage{array}
\usepackage{enumitem}
\usepackage{microtype}
\usepackage{hyperref}

\geometry{left=2.6cm,right=2.6cm,top=2.5cm,bottom=2.5cm}
\setlength{\parindent}{2em}
\setlength{\parskip}{0.35em}
\renewcommand{\arraystretch}{1.35}
\setlist{nosep}
\hypersetup{hidelinks}

\title{方法与实验部分}
\author{}
\date{}

\begin{document}

\maketitle

\section{3 队列感知动态路径组织模型（AA）}

大型多线换乘站应急疏散的难点不只在于寻找几何距离较短的路径，而在于多线路站台客流、列车到达客流和换乘客流会在楼梯、扶梯、换乘通道、闸机和出口等共享设施处持续汇聚。若路径组织只依据当前节点状态或静态边权进行判断，某一批次在决策时看似可行的路径，可能在其真正到达关键设施时已经面对由前序批次形成的排队和下游空间回溢。为描述这一过程，本文提出队列感知动态路径组织模型 AdaptiveQueueAwareAStar，以下简称 AA。

AA 不是单一 A* 启发式权重的修改，而是由路径决策层和设施服务执行层耦合形成的疏散组织模型。路径决策层在批次进入选择阶段时预测其到达共享设施时的等待状态，并据此搜索通向出口的可执行路径；设施服务执行层按照统一的设施容量、整数服务能力和下游空间接收约束执行已经给出的移动请求。前者回答“乘客应被组织到哪条路径”，后者回答“这条路径在设施容量约束下能否、何时被实际通过”。两层相互连接，使路径选择不脱离后续排队和通行过程。

\subsection{3.1 AA 模型总体结构}

AA 的输入包括站内疏散网络、来源批次、候选出口、共享设施容量、当前设施队列、下游空间占用状态以及已经承诺的未来到达事件。模型输出不是一条全站统一的固定疏散路线，而是面向不同来源批次的路径组织结果，包括批次下一步移动方向、完整候选路径、预计到达关键设施的时间、预计等待时间以及相应的未来到达事件登记。

设站内疏散系统在时刻 \(t\) 的状态为
\[
\Omega(t)=\left(G,B(t),Q(t),N(t),A(t)\right),
\]
其中，\(G\) 为站内有向网络，\(B(t)\) 为进入决策阶段的客流批次集合，\(Q(t)\) 为共享设施队列状态，\(N(t)\) 为空间节点占用状态，\(A(t)\) 为已经承诺的未来设施到达事件。AA 在每个离散时间步更新 \(\Omega(t)\)，并对可决策批次依次进行路径搜索和事件登记。由于前一批次登记的未来到达事件会改变后一批次看到的设施负荷，AA 能够显式刻画多源客流在时间上的先后竞争关系。

模型由四个子模块组成。第一，站内疏散网络与来源批次模块，用于描述五线换乘站内部空间、设施连接和不同来源客流进入疏散系统的时间结构。第二，到达时刻设施负荷预测模块，用于预测候选批次到达楼扶梯、通道、闸机和出口时的队列状态。第三，时间依赖 AA 路径决策模块，用预计行走时间、设施等待和空间回溢等待构造路径代价，并对不同出口路径进行搜索。第四，共享容量执行模块，用同一套设施服务能力执行 AA 和对照方法给出的移动请求，保证比较差异来自路径组织逻辑，而不是来自执行规则差异。

\subsection{3.2 站内疏散网络与来源批次}

\subsubsection{3.2.1 站内设施网络}

将龙阳路站公共区抽象为有向图
\[
G=(V,E),
\]
其中，节点集合 \(V\) 表示线路站台、站厅、换乘连接节点、楼梯、扶梯、闸机、出口前区域和地面出口，边集合 \(E\) 表示相邻空间或设施之间允许乘客移动的方向性连接。每条边 \(e=(i,j)\) 包含长度 \(l_e\)、有效宽度 \(w_e\)、行进方向、设施类型和容量属性；每个节点 \(v\) 包含面积 \(A_v\)、当前人数 \(N_v(t)\)、节点类型和空间接收上限。

大型换乘站的设施容量不能简单理解为每条边各自独立的通行能力。同一楼扶梯、同一换乘通道或同一闸机组可能同时承接来自多个站台、多个线路和多个换乘方向的乘客。为刻画这种竞争关系，模型建立边到共享设施资源的映射。若边 \(e\) 需要消耗共享设施 \(r\) 的通行能力，则记为
\[
\rho(e)=r.
\]
映射到同一资源 \(r\) 的多条边共同竞争该资源的服务能力；若边只表示普通空间移动而不消耗独立瓶颈资源，则其资源映射为空。该处理使换乘客流与出站客流、列车到达客流在同一设施上的叠加能够被直接记录，而不是被拆分为若干互不影响的边容量。

\subsubsection{3.2.2 来源批次与换乘客流}

设客流来源组集合为 \(\mathcal{G}\)。来源组用于区分不同线路站台、列车到达、站厅滞留和跨线换乘客流。对龙阳路站这类大型五线换乘站而言，换乘客流不是独立于整体疏散之外的对象，而是改变全站设施负荷分布的重要来源。其路径往往跨越多个公共区和垂向设施，并在疏散过程中与站台释放客流和列车到达客流共同占用关键通道、楼扶梯、闸机和出口。

客流以自然到达批次表示。一个批次记为
\[
b=(g,n_b,x_b,t_b),
\]
其中，\(g\) 为来源组，\(n_b\) 为批次人数，\(x_b\) 为当前节点，\(t_b\) 为批次进入路径决策的时间。采用批次作为路径决策单元，可以保留列车到达、站台释放和换乘流进入公共区的时间结构，也能使后续搜索显式记录不同批次到达同一共享设施的先后顺序。对已经进入边内移动或已经获得设施服务能力的批次，模型不在途中瞬时改变其移动方向；只有仍处于明确选择阶段的批次才参与 AA 的路径重评估。

\subsubsection{3.2.3 共享设施服务队列}

对具有服务能力限制的设施 \(r\)，设其服务率为 \(\mu_r\)，单位为人每秒，仿真步长为 \(\Delta t\)。在时间步 \(t\) 内，设施 \(r\) 可释放的整数服务能力为
\[
C_r(t)=\left\lfloor \mu_r\Delta t+\xi_r(t)\right\rfloor ,
\]
其中，\(\xi_r(t)\) 为上一时间步未形成完整整数人的容量余量。设 \(A_r(t)\) 为时间步 \(t\) 内请求进入设施 \(r\) 的人数，\(Q_r(t)\) 为该设施入口等待队列，则队列递推为
\[
Q_r(t+\Delta t)=\mathrm{max}\left(Q_r(t)+A_r(t)-C_r(t),0\right).
\]

这一表达式体现了共享设施的服务属性。楼梯、扶梯、通道和闸机不是路径图上的普通连接线，而是具有服务释放过程的瓶颈资源。若多个来源组在同一时间步到达同一资源，它们共同竞争同一 \(C_r(t)\)，未获得服务能力的人数停留在上游等待。由此，路径搜索给出的方向不会被直接等同于实际通行，实际通行必须通过容量执行层完成。

\subsubsection{3.2.4 空间接收与物理行走时间}

设施服务能力还受到下游空间接收条件限制。设节点 \(v\) 的空间接收上限为 \(\bar N_v\)，当前占用人数为 \(N_v(t)\)，当前时间步已经被其他移动请求预留进入的人数为 \(B_v(t)\)，则新进入人数 \(y_v(t)\) 需要满足
\[
N_v(t)+B_v(t)+y_v(t)\leq \bar N_v .
\]
当下游节点接收不足时，即使上游设施仍有名义服务能力，实际移动也可能被阻滞。这一约束用于描述高负荷换乘站中常见的回溢现象，即队列不仅在闸机或楼梯入口形成，也可能从下游通道、站厅或出口前区域向上游扩散。

对于普通移动边，行走时间由边长和速度决定。当前实现采用速度-密度关系计算边内通行时间：
\[
v(k)=
\begin{cases}
v_f, & k\leq k_f,\\
\mathrm{max}(v_f-ak,0), & k_f<k<k_j,\\
0, & k\geq k_j .
\end{cases}
\]
其中，\(k\) 为边内客流密度，\(v_f\) 为自由流速度，\(k_f\) 为自由流密度阈值，\(a\) 为速度下降斜率，\(k_j\) 为拥挤密度上限。本文采用的 \(v_f=1.427\ \mathrm{m/s}\)、\(k_f=0.2\ \mathrm{p/m^2}\)、\(a=0.3549\) 和 \(k_j=4.0\ \mathrm{p/m^2}\) 来自对改进 A* 疏散仿真文献中速度-密度关系的复现，不作为龙阳路站实测速度。该关系用于计算物理行走时间，但本文的适用边界分析不把速度-密度阈值作为核心变量。

\subsection{3.3 到达时刻设施负荷预测}

AA 的关键在于不只读取当前设施队列，而是预测候选批次真正到达共享设施时的队列状态。低负荷条件下，当前队列与到达时队列之间的差异可能较小；高负荷条件下，多线路站台客流、列车到达客流和换乘客流会在不同时间先后到达同一设施，当前时刻尚未出现的排队可能在候选批次到达时已经形成。若忽略这一时间差，路径组织容易低估共享设施未来负荷。

对共享设施 \(r\)，定义从当前时刻 \(t\) 开始已经承诺的未来到达事件集合为
\[
\mathcal{A}_r(t)=\{(\eta_j,n_j)\},
\]
其中，\(\eta_j\) 为已确定批次预计进入设施 \(r\) 的时刻，\(n_j\) 为对应人数。对候选批次 \(b\)，若其预计在时刻 \(\eta\) 到达设施 \(r\)，则设施到达队列估计为
\[
\widehat Q_r(\eta)=\mathcal{Q}\left(Q_r(t),\mathcal{A}_r(t),\mu_r,t,\eta\right).
\]
队列演化算子 \(\mathcal{Q}\) 在 \(t\) 到 \(\eta\) 的区间内释放服务能力，并在已登记到达事件发生时加入相应人数。若 \(\mu_r>0\)，候选批次在设施 \(r\) 处的预计等待时间为
\[
W_r(\eta)=\frac{\widehat Q_r(\eta)}{\mu_r}.
\]

未来到达事件包括两类。第一类是已经进入执行过程、到达时间可以由当前移动状态推算的批次；第二类是同一决策轮中已经完成 AA 搜索、但尚未实际到达设施的前序批次。当前序批次路径确定后，其预计进入共享设施的时间和人数立即登记为到达事件，使后续批次能够在搜索时看到这一已经承诺的设施负荷。由此，AA 在同一时间步内部也能形成先后顺序，而不是让所有批次基于同一份未更新的空白预测状态独立决策。

该预测模块是确定性队列递推，不是机器学习预测器。它服务于本文的路径组织问题：把“当前是否拥堵”的判断推进到“候选批次到达时是否拥堵”的判断。对于换乘客流占比较高或列车到达高度重叠的场景，这一差别直接影响候选出口和关键设施路径的排序。

\subsection{3.4 AA 路径代价与决策模型}

对批次 \(b\) 的候选路径 \(P=(e_1,e_2,\ldots,e_m)\)，设第 \(i\) 条边的预计进入时刻为 \(\eta_i\)，对应共享设施为 \(r_i=\rho(e_i)\)。AA 的路径代价定义为
\[
J_b(P)=\sum_{i=1}^{m}\left[\tau_{e_i}(\eta_i)+W_{r_i}(\eta_i)+S_{r_i}(n_b)+H_{e_i}(\eta_i)\right],
\]
其中，\(\tau_{e_i}\) 为物理行走时间，\(W_{r_i}\) 为到达设施时的预计排队等待，\(S_{r_i}(n_b)\) 为当前批次自身占用服务能力形成的平均服务时间，\(H_{e_i}\) 为下游空间接收不足或回溢造成的预计等待。若边不对应共享设施，则相关服务等待项取 0。

路径代价按预计到达时刻递推。前一段路径上的行走和等待会改变后一段路径的到达时间，后一段设施的等待又取决于该新的到达时间。因此，AA 不是静态最短路，也不是只把当前密度写入边权的即时重算，而是面向设施到达负荷的时间依赖路径搜索。

路径搜索标签定义为
\[
L=(v,\eta,g,\pi),
\]
其中，\(v\) 为当前节点，\(\eta\) 为预计到达当前节点的时刻，\(g\) 为累计路径代价，\(\pi\) 为前驱标签。标签从节点 \(u\) 沿边 \(e=(u,v)\) 扩展时，先计算物理行走时间，再根据预计进入设施时刻查询共享设施队列和下游空间接收状态，得到新的到达时刻与累计代价。两个标签即使到达同一节点，只要到达时刻不同，后续设施队列状态就可能不同，因此不能简单按节点编号合并为一个静态标签。静态自由流时间仅作为启发式下界使用，不替代实际路径代价。

AA 不使用固定 \(K\) 条候选路径，也不缓存动态完整路径。其缓存只用于静态自由流启发式下界、共享资源映射和到达事件索引。这样可以避免同一节点的不同批次在设施预测状态已经变化后仍复用第一批次的路径结果。

对仍处于选择阶段的批次，AA 允许在满足稳定性条件时重新评估路径。设保留当前路径的剩余代价为 \(C_{\mathrm{stay}}\)，候选切换路径代价为 \(C_{\mathrm{switch}}\)，相对收益为
\[
R=\frac{C_{\mathrm{stay}}-C_{\mathrm{switch}}}{C_{\mathrm{stay}}}.
\]
仅当候选路径有效、下一步移动方向不同、批次仍处于可重评估阶段、\(C_{\mathrm{switch}}<C_{\mathrm{stay}}\)，且 \(R\) 超过防摆动阈值时，路径切换才被接受。防摆动阈值是算法稳定性设置，不表述为行人行为参数。

\subsection{3.5 疏散效果评价指标}

评价指标服务于整体疏散机制解释，而不是只给出某一方法的平均时间优劣。本文从四个层面评价路径组织效果。

第一，整体疏散效率。采用 \(T_{95}\)、\(T_{100}\)、平均疏散时间、累计停滞人秒、总移动距离和有效疏散速度描述全站清空过程。其中，\(T_{95}\) 用于刻画系统进入尾部清空阶段前的效率，\(T_{100}\) 用于描述最后一批乘客离站时间，累计停滞人秒用于衡量乘客在排队、阻塞和等待状态中的暴露程度。

第二，关键设施瓶颈。对楼扶梯、换乘通道、闸机和出口前区域记录通过人数、峰值队列、队列人秒、首次排队时间、最后排队时间和利用率。若一种方法降低系统尾部时间但增加局部设施队列，该变化必须在该层指标中呈现。

第三，来源组和换乘客流分解。按线路站台、列车到达、站厅和换乘来源统计出口分配、设施通过量和清空时间，用于解释不同来源客流在哪些共享设施上汇聚，以及换乘客流对全站尾部疏散的贡献。

第四，跨模型一致性。将宏观网络模型的路径组织结果导入 Pathfinder 微观仿真，比较疏散曲线、拥堵空间分布、出口利用和乘客拥堵时间。Pathfinder 不被视为现场观测真值，而用于检验宏观路径组织结果在微观行人运动尺度上是否呈现相近趋势。

\section{4 AA 模型求解与仿真实现}

第 3 章给出了 AA 的模型组成，本章说明该模型如何在离散仿真中被求解和执行。这里的“求解”不是一次性求出全站固定路线，而是在疏散过程中随客流批次到达持续更新路径决策、未来设施负荷和共享容量执行结果。

\subsection{4.1 从模型状态到可执行搜索标签}

在仿真时刻 \(t\)，系统首先读取当前状态 \(\Omega(t)\)，包括节点占用、设施队列、在途批次、可决策批次和已登记的未来到达事件。对每一个可决策批次，AA 将第 3 章中的路径代价转化为时间依赖 A* 标签搜索。搜索入口对应批次当前节点，目标为可用出口集合，搜索过程中每一次边扩展都通过共享执行层提供的回调函数读取物理行走时间、设施服务能力和空间接收状态。

标签扩展过程包含三个步骤。第一，根据边长、设施类型和速度-密度关系计算从 \(u\) 到 \(v\) 的物理行走时间 \(\tau_e\)。第二，若该边对应共享设施 \(r\)，则根据预计进入时刻查询 \(\widehat Q_r(\eta)\)，得到预计设施等待 \(W_r(\eta)\) 和当前批次自身服务时间 \(S_r(n_b)\)。第三，检查下游节点或边空间的接收条件，估计空间阻塞等待 \(H_e(\eta)\)。三项结果共同更新标签到达时刻和累计代价。

AA 的启发式函数只提供静态自由流条件下的时间下界，用于减少搜索空间；它不代表真实路径代价。真实代价仍由物理行走时间、到达时刻设施等待、当前批次服务时间和空间接收等待共同决定。该设置保证了搜索效率，同时保留高负荷换乘场景中“到达时才出现的瓶颈”对路径选择的影响。

\subsection{4.2 AA 搜索与到达事件登记流程}

AA 以离散时间步推进，每一时间步包含状态更新、批次决策、事件登记、容量执行和指标记录五个环节。

\begin{enumerate}
    \item 更新上一时间步已经完成移动的乘客，将其释放到目标节点、设施队列或出口统计中。
    \item 根据站台释放、列车到达和换乘来源规则生成本时间步进入选择阶段的自然到达批次。
    \item 按到达时间、来源组、当前节点和批次编号对可决策批次排序，保证同一轮决策具有确定性。
    \item 对排序后的每个批次执行时间依赖 A* 搜索，计算候选路径的预计行走时间、设施等待、当前批次服务时间和空间阻塞等待。
    \item 路径确定后，将该批次预计进入共享设施的时刻和人数登记为未来到达事件，使同一轮中的后续批次能够看到已经承诺的设施负荷。
    \item 对仍处于选择阶段且满足稳定性条件的批次进行路径重评估，已经进入在途移动或已经获得设施服务能力的批次继续执行已承诺移动。
    \item 汇总所有上游节点对同一共享设施的移动请求，交由共享容量执行层分配本时间步实际通过人数。
    \item 记录全站疏散、关键设施队列、来源组位置、出口分配、路径重评估和计算耗时等指标。
\end{enumerate}

到达事件登记是 AA 与普通动态 A* 的关键区别之一。若前序批次已经选择经过某一楼扶梯或闸机，即使其尚未物理到达该设施，其预计到达时间和人数也会被写入事件集合。后序批次搜索同一设施时，预测队列中已包含这部分承诺负荷。因此，AA 不是简单让所有批次各自找一条短路，而是在同一仿真轮内部建立批次之间的设施竞争关系。

\subsection{4.3 共享容量执行与守恒约束}

路径搜索只产生移动意图，实际通行由共享容量执行层完成。对同一设施 \(r\)，无论移动请求来自哪个来源组、哪条线路或哪条边，只要它们映射到同一资源，就共同使用该资源在当前时间步的整数服务能力 \(C_r(t)\)。当需求超过服务能力时，未被服务的人数保留在上游队列，并在后续时间步继续参与容量竞争。

共享容量执行层同时检查下游空间接收条件。若下游节点达到接收上限，即使设施本身仍有剩余服务能力，相关移动也不能被直接释放到下游空间。这一处理避免路径搜索与物理执行不一致：路径上可达不代表该时间步可通过，设施有能力不代表下游空间能够接收。

该执行层对 AA 和 ImprovedAStar 完全相同。这样设置的目的不是增加一个额外模型，而是把比较对象限定在路径组织逻辑本身。若两种方法的客流输入、站内网络、设施服务率、速度-密度关系和空间接收规则不同，则无法判断结果差异来自路径组织，还是来自基础仿真条件。统一执行层使后续实验能够把差异归因到“是否预测到达时刻设施负荷”这一核心机制。

\subsection{4.4 计算对照方法}

ImprovedAStar 在本文中仅作为计算对照方法，不代表龙阳路站运营方制定或发布的官方疏散预案。该方法在同一站内网络、同一客流输入和同一执行层下进行路径搜索，但不使用未来设施到达事件集合，不预测候选批次到达设施时的排队状态，也不执行基于收益阈值的路径重评估。它反映的是一种传统改进 A* 路径组织基准，用于检验 AA 的到达负荷预测和共享容量耦合是否能改善整体疏散过程。

对照实验采用共同终止规则和共同指标体系。两种方法均以目标客流全部到达出口为终止条件，均记录整体疏散时间、尾部清空时间、停滞人秒、关键设施队列、来源组清空时间、出口利用和计算耗时。若 AA 缩短尾部时间但增加移动距离或计算耗时，结果部分应同时报告该权衡；若 ImprovedAStar 在低负荷下表现接近 AA，也应作为方法适用边界的一部分呈现。

\section{5 案例研究与实验设计}

本文以龙阳路站为案例开展实验。龙阳路站为 2 号线、7 号线、16 号线、18 号线及磁浮线的五线换乘站，站内换乘关系复杂，具有多线路站台、换乘通道、楼扶梯、闸机和多出口共同作用的典型特征。模型根据 CAD 图纸和仿真配置建立站台、站厅、换乘通道、楼扶梯、闸机和出口网络；设施宽度、闸机数量和出口宽度在 CAD 核对完成后进入最终参数表。未能由公开资料或 CAD 直接确认的参数，不写作实测值，而在参数表中标记为模型设定值或文献取值。

\subsection{5.1 案例车站与疏散场景}

实验设置低负荷常规应急场景和高负荷满载列车叠加场景。低负荷常规应急场景包含各线路站台等待、站厅和换乘基础客流，输入总人数为 2187 人。该场景用于检验在设施竞争较弱时，AA 是否保持稳定，是否因过度预测引入不必要绕行或出口偏置。

高负荷满载列车叠加场景在基础客流上加入多线路列车到达客流，输入总人数为 17905 人。该场景用于放大多线路站台客流、列车客流和换乘客流在共享设施上的汇聚效应，检验到达时刻设施负荷预测能否降低整体等待暴露和尾部滞留。

两类场景均比较 ImprovedAStar 与 AA。两种方法使用相同的站内网络、客流输入、设施服务能力、速度-密度关系、空间接收约束和终止规则。最终结果冻结前，正文不写入百分比改善或最终优劣结论。

\begin{table}[htbp]
\centering
\caption{低负荷与高负荷疏散场景设置}
\small
\begin{tabular}{p{0.20\textwidth}p{0.16\textwidth}p{0.25\textwidth}p{0.27\textwidth}}
\toprule
场景 & 输入人数 & 主要客流构成 & 实验目的 \\
\midrule
低负荷常规应急 & 2187 & 站台等待、站厅滞留和基础换乘客流 & 检验设施竞争较弱时 AA 的稳定性和低负荷适用性 \\
高负荷满载列车叠加 & 17905 & 基础客流叠加多线路列车到达客流 & 检验多来源客流同步汇聚时 AA 对等待暴露和尾部滞留的影响 \\
\bottomrule
\end{tabular}
\end{table}

\subsection{5.2 AA 到达负荷预测结果分析}

案例研究首先验证 AA 的中间预测模块，而不是直接进入总体时间比较。该实验回答三个问题：第一，AA 是否把前序批次已经承诺的到达事件计入后续批次的设施等待；第二，预计等待是否改变候选路径排序；第三，预测得到的关键设施压力是否能在执行层实际队列中得到对应体现。

实验选取换乘通道、楼扶梯、闸机和出口前区域中的关键共享设施作为分析对象。对每个设施记录实际队列、未来确定到达人数、预测等待时间、实际通过人数和未获得服务能力人数。对代表性来源组，进一步分解其候选路径代价，包括物理行走时间、设施等待时间、空间回溢等待和总路径代价。结果冻结后，本节给出关键设施队列曲线和代表性批次路径代价分解图，用于说明到达负荷预测如何改变路径组织判断。

\begin{table}[htbp]
\centering
\caption{AA 到达负荷预测分析指标}
\small
\begin{tabular}{p{0.20\textwidth}p{0.34\textwidth}p{0.34\textwidth}}
\toprule
分析对象 & 指标 & 作用 \\
\midrule
关键楼扶梯与换乘通道 & 预测到达人数、实际队列、峰值等待、最后排队时间 & 判断多来源客流是否在垂向设施和换乘连接处集中汇聚 \\
闸机与出口前区域 & 通过人数、队列人秒、出口分配、来源组构成 & 判断出口方向选择是否造成局部服务压力 \\
代表性来源组批次 & 路径代价分解、候选路径排序、是否触发路径重评估 & 解释路径改变来自行走距离、设施等待还是空间回溢 \\
\bottomrule
\end{tabular}
\end{table}

\subsection{5.3 疏散路径组织方案性能分析}

整体性能比较 \(T_{95}\)、\(T_{100}\)、平均疏散时间、累计停滞人秒、平均移动时间、总移动距离和计算耗时。低负荷场景重点观察 AA 是否保持稳定；高负荷场景重点观察尾部清空、停滞等待和关键设施排队。若某方法降低等待但增加移动距离，或者缩短系统尾部但增加局部设施队列，正文需要同时报告这些权衡。

\begin{table}[htbp]
\centering
\caption{低负荷与高负荷场景下的整体疏散性能记录表}
\small
\begin{tabular}{p{0.18\textwidth}p{0.16\textwidth}p{0.10\textwidth}p{0.12\textwidth}p{0.12\textwidth}p{0.14\textwidth}p{0.12\textwidth}}
\toprule
场景 & 方法 & 人数 & \(T_{95}\) & \(T_{100}\) & 累计停滞人秒 & 总移动距离 \\
\midrule
低负荷常规应急 & ImprovedAStar & 2187 & 待冻结 & 待冻结 & 待冻结 & 待冻结 \\
低负荷常规应急 & AA & 2187 & 待冻结 & 待冻结 & 待冻结 & 待冻结 \\
高负荷满载列车叠加 & ImprovedAStar & 17905 & 待冻结 & 待冻结 & 待冻结 & 待冻结 \\
高负荷满载列车叠加 & AA & 17905 & 待冻结 & 待冻结 & 待冻结 & 待冻结 \\
\bottomrule
\end{tabular}
\end{table}

在整体指标之后，进一步分析关键设施队列和来源组分解。该分析用于解释整体差异的物理原因，而不是作为附加结果。若高负荷场景的尾部清空时间发生变化，需要追踪这种变化是否来自某些换乘通道、楼扶梯或闸机的持续排队减少；若系统尾部时间缩短但某一线路或某一区域排队增加，也需要在该层结果中说明局部代价。

来源组分解重点关注换乘客流与站台、列车来源客流在共享设施上的叠加关系。对每一类来源组统计其主要路径链、关键设施通过量、出口选择和清空时间。通过该分析回答三个问题：换乘客流是否是某些瓶颈形成的重要组成部分；AA 是否改变了换乘客流到达关键设施的时间分布；这种变化是否进一步影响全站整体疏散效率。

\subsection{5.4 综合效果、适用边界与 Pathfinder 对照}

综合效果分析首先回到大型换乘站的空间机制。对高负荷场景，按照关键设施队列人秒和峰值队列识别主要瓶颈，再将瓶颈设施的客流来源分解为不同线路站台、列车到达和换乘来源。若同一瓶颈由多个来源共同构成，说明该设施是全站疏散组织中的共享限制环节；若某一方法只是把等待从一个设施转移到另一个设施，则应在路径链和设施队列结果中明确呈现。

本文保留适用边界分析，但不把速度-密度阈值作为主实验变量。速度-密度关系属于行走时间模型参数，对本文核心问题的解释力有限。更能反映多线换乘站疏散组织边界的变量包括三类：换乘客流比例、列车到达重叠程度和关键共享设施服务能力。换乘客流比例用于检验方法是否主要在换乘客流显著参与设施竞争时发挥作用；列车到达重叠程度用于检验多来源客流同时释放对设施排队和尾部清空的影响；关键共享设施服务能力用于检验方法在瓶颈强弱变化下的适用范围。

\begin{table}[htbp]
\centering
\caption{适用边界分析的变量设置}
\small
\begin{tabular}{p{0.22\textwidth}p{0.28\textwidth}p{0.38\textwidth}}
\toprule
变量类别 & 设置方式 & 分析目的 \\
\midrule
换乘客流比例 & 在总客流规模基本一致的前提下调整换乘来源占比 & 判断换乘客流汇聚是否是 AA 产生差异的重要场景条件 \\
列车到达重叠程度 & 调整多线路列车客流释放的时间间隔 & 判断同步到达和错峰到达对共享设施排队的影响 \\
关键共享设施服务能力 & 对主要楼扶梯、换乘通道或闸机服务率进行扰动 & 判断路径组织效果对瓶颈强度的依赖程度 \\
\bottomrule
\end{tabular}
\end{table}

Pathfinder 对照用于检验宏观网络模型的路径组织结果在微观行人仿真中是否呈现一致趋势。本文将宏观模型导出的来源组、路径链和路径人数转换为 Pathfinder 中的行人组和行为路径，并在相同车站几何和出口设置下运行微观仿真。对照指标包括累计疏散曲线、尾部清空时间、拥堵时间、关键设施附近空间拥挤分布和出口利用。

该对照不把 Pathfinder 视为真实车站观测，也不以 Pathfinder 结果替代宏观模型结果。其作用是检查：宏观模型识别出的瓶颈位置是否在微观轨迹中出现相近拥堵；路径组织引起的出口和设施使用变化是否能在微观模型中复现；若两者存在差异，则进一步从导航网格、个体避让、门选择规则和容量映射差异解释。

结果冻结后，案例研究按照以下顺序写作：首先报告 AA 到达负荷预测结果，证明中间预测模块参与并改变路径评价；其次比较低负荷和高负荷场景下两种路径组织方法的整体疏散效率；再次通过关键设施和来源组分解解释整体差异的站内机制，重点说明换乘客流与其他来源客流在共享设施上的汇聚关系；随后通过适用边界分析说明方法在不同换乘比例、到达重叠和瓶颈强度下的有效范围；最后用 Pathfinder 微观仿真对宏观疏散趋势进行交叉验证。

\end{document}
